\chapter{Modelagem de Casos de Uso}
\label{cap_modelagem_casos_uso}

A modelagem de casos de uso na notação UML 2.5 (\textit{Unified Modeling Language}) estabelece a ponte conceitual entre a especificação declarativa de requisitos e a arquitetura dinâmica de interação do software, permitindo mapear como os atores realizam suas metas no sistema \cite{pressman2021,booch2005}.

% ===========================================================================
\section{Diagrama Geral de Casos de Uso}
\label{sec_diagrama_casos_uso}
% ===========================================================================

A \autoref{fig_usecase_requisitos} ilustra a delimitação dos casos de uso do SIGEA-GTT e sua articulação com os três atores do sistema (\texttt{ADMIN}, \texttt{PROFESSOR} e \texttt{STUDENT}), derivados diretamente dos 25 Requisitos Funcionais.

\begin{figure}[p]
\centering
\caption{Diagrama Geral de Casos de Uso do SIGEA-GTT (UML 2.5).}
\label{fig_usecase_requisitos}
\resizebox{0.92\textwidth}{!}{%
\begin{tikzpicture}[
  >=Stealth,
  uc/.style={
    ellipse, draw=blue!70!black, fill=blue!5, thick,
    font=\tiny, align=center, minimum width=3.2cm, minimum height=0.75cm
  },
  bdr/.style={
    rectangle, draw=blue!40, dashed, thick, rounded corners=6pt,
    fill=blue!2
  },
  act/.style={font=\small\bfseries, draw=blue!80!black, fill=white, rounded corners=2pt, inner sep=4pt}
]

% Sistema Boundary
\node[bdr, minimum width=9.5cm, minimum height=15.5cm]
     (sys) at (5, -6) {};
\node[above, font=\small\bfseries, blue!80!black] at (sys.north) {SIGEA-GTT -- Sistema de Gestão de Eventos Adversos};

% Atores
\node[act] (admin) at (-2,  -1.5) {ADMIN};
\node[act] (prof)  at (-2,  -7.5) {PROFESSOR};
\node[act] (aluno) at (12,  -8.5) {STUDENT};

% Casos de uso comuns
\node[uc] (login)   at (5,   0.5) {UC01: Autenticar no Sistema\\(RF01, RF02)};
\node[uc] (senha)   at (5,  -0.8) {UC02: Alterar Senha Pessoal\\(RF03)};
\node[uc] (catalogo)at (5,  -2.1) {UC03: Consultar Catálogo GTT\\e Guia MERP (RF08, RF09)};

% Casos de uso Admin
\node[uc] (uc_users)   at (3, -3.8) {UC04: Gerenciar Usuários\\(RF04, RF05)};
\node[uc] (uc_modgtt)  at (3, -5.2) {UC05: CRUD Módulos GTT\\(RF06)};
\node[uc] (uc_gat)     at (3, -6.6) {UC06: CRUD 53 Gatilhos IHI\\(RF07)};
\node[uc] (uc_turmas)  at (3, -8.0) {UC07: Gerenciar Turmas\\(RF10, RF11, RF12, RF13)};
\node[uc] (uc_dash_a)  at (3, -9.5) {UC08: Monitorar Dashboard\\Geral (RF25)};

% Casos de uso Professor
\node[uc] (uc_atv)     at (7, -3.8) {UC09: Elaborar Atividade e\\Prontuário JSONB (RF14, RF15)};
\node[uc] (uc_corr)    at (7, -5.2) {UC10: Avaliar Submissão e\\Emitir Parecer (RF24)};
\node[uc] (uc_dash_p)  at (7, -6.6) {UC11: Analisar Indicadores da\\Turma (RF25)};

% Casos de uso Aluno
\node[uc] (uc_cron)    at (7, -8.2) {UC12: Auditar Prontuário sob\\Cronômetro 20 min (RF16)};
\node[uc] (uc_gatid)   at (7, -9.6) {UC13: Identificar Gatilhos\\e Danos MERP (RF17)};
\node[uc] (uc_ferrq)   at (7,-11.0) {UC14: Aplicar Ferramentas da\\Qualidade (RF18 a RF22)};
\node[uc] (uc_sub)     at (7,-12.4) {UC15: Submeter Resolução\\Definitiva (RF23)};

% Relacionamentos Admin
\draw[->] (admin) -- (login);
\draw[->] (admin) -- (uc_users);
\draw[->] (admin) -- (uc_modgtt);
\draw[->] (admin) -- (uc_gat);
\draw[->] (admin) -- (uc_turmas);
\draw[->] (admin) -- (uc_dash_a);

% Relacionamentos Professor
\draw[->] (prof) -- (login);
\draw[->] (prof) -- (catalogo);
\draw[->] (prof) -- (uc_atv);
\draw[->] (prof) -- (uc_corr);
\draw[->] (prof) -- (uc_dash_p);

% Relacionamentos Aluno
\draw[->] (aluno) -- (login);
\draw[->] (aluno) -- (catalogo);
\draw[->] (aluno) -- (uc_cron);
\draw[->] (aluno) -- (uc_gatid);
\draw[->] (aluno) -- (uc_ferrq);
\draw[->] (aluno) -- (uc_sub);

% Include
\draw[->, dashed] (login) -- node[right, font=\tiny]{include} (senha);
\draw[->, dashed] (uc_cron) -- (uc_gatid);
\draw[->, dashed] (uc_gatid) -- (uc_ferrq);
\draw[->, dashed] (uc_ferrq) -- (uc_sub);

\end{tikzpicture}%
}
\fonte{Elaborado pelo autor com TikZ (2026).}
\end{figure}

% ===========================================================================
\section{Especificação dos Casos de Uso Críticos}
\label{sec_especificacao_casos_uso}
% ===========================================================================

Para assegurar a completude e testabilidade da especificação, os casos de uso de maior complexidade comportamental são detalhados a seguir através do formato textual padronizado de Cockburn \cite{pressman2021}.

% ---------------------------------------------------------------------------
\subsection{UC01: Autenticar e Redefinir Senha Inicial}
\label{sub_uc01}
% ---------------------------------------------------------------------------
\begin{itemize}
    \item \textbf{Atores Primários}: Qualquer usuário cadastrado (\texttt{ADMIN}, \texttt{PROFESSOR}, \texttt{STUDENT});
    \item \textbf{Requisitos Vinculados}: \texttt{RF01}, \texttt{RF02}, \texttt{RF03}, \texttt{RNF03};
    \item \textbf{Pré-condição}: O usuário possui conta ativa no domínio \texttt{@academico.ufs.br};
    \item \textbf{Fluxo Principal}:
    \begin{enumerate}
        \item O usuário acessa a página de autenticação e informa seu e-mail institucional e senha (provisória ou cadastrada);
        \item O sistema valida as credenciais contra o hash BCrypt persistido;
        \item O sistema emite o token JWT assinado;
        \item Se a flag \texttt{primeiroAcesso} estiver ativa, o sistema redireciona imediatamente para a tela de troca obrigatória de senha;
        \item O usuário digita a nova senha desejada e sua confirmação;
        \item O sistema valida a complexidade da senha, atualiza o hash no banco, desativa a flag \texttt{primeiroAcesso} e redireciona o usuário para seu painel principal.
    \end{enumerate}
    \item \textbf{Fluxos de Exceção}:
    \begin{itemize}
        \item \textbf{E-mail fora do domínio \texttt{@academico.ufs.br}}: O sistema bloqueia o envio com mensagem de validação imediata;
        \item \textbf{Credenciais Inválidas}: O sistema emite aviso de erro HTTP 401 e sugere contatar o gestor acadêmico.
    \end{itemize}
\end{itemize}

% ---------------------------------------------------------------------------
\subsection{UC07: Gerenciar Turmas Acadêmicas}
\label{sub_uc07}
% ---------------------------------------------------------------------------
\begin{itemize}
    \item \textbf{Atores Primários}: Administrador (\texttt{ADMIN});
    \item \textbf{Requisitos Vinculados}: \texttt{RF10}, \texttt{RF11}, \texttt{RF12}, \texttt{RF13}, \texttt{RNF05}, \texttt{RNF08};
    \item \textbf{Pré-condição}: O administrador está autenticado; existem docentes cadastrados no sistema;
    \item \textbf{Fluxo Principal}:
    \begin{enumerate}
        \item O administrador solicita a criação de uma nova turma acadêmica;
        \item O administrador preenche a matéria, turma e período (gerando o padrão obrigatório \texttt{"Matéria - Turma - Período"}) e seleciona o docente titular;
        \item O sistema cria a turma associando atomicamente o professor titular;
        \item O administrador seleciona discentes cadastrados para matrícula na turma;
        \item O sistema valida que todos os discentes possuem perfil \texttt{STUDENT} e efetiva os vínculos.
    \end{enumerate}
    \item \textbf{Fluxo Alternativo (Encerramento da Turma)}:
    \begin{enumerate}
        \item O administrador solicita o encerramento da turma;
        \item O sistema verifica se todas as submissões foram avaliadas pelo docente titular;
        \item Em caso positivo, o sistema exibe diálogo modal de confirmação em duas etapas;
        \item Confirmada a ação, a turma é marcada como \texttt{ENCERRADA}.
    \end{enumerate}
    \item \textbf{Fluxo de Exceção (Pendências de Correção)}:
    \begin{itemize}
        \item Existindo submissões sem nota ou parecer pedagógico, o sistema bloqueia o encerramento e exibe alerta detalhado impedindo a operação (\texttt{RF13}).
    \end{itemize}
\end{itemize}

% ---------------------------------------------------------------------------
\subsection{UC12/13/14: Auditar Prontuário Simulado sob Cronômetro e Aplicar Ferramentas}
\label{sub_uc12_14}
% ---------------------------------------------------------------------------
\begin{itemize}
    \item \textbf{Atores Primários}: Discente (\texttt{STUDENT});
    \item \textbf{Requisitos Vinculados}: \texttt{RF15}, \texttt{RF16}, \texttt{RF17}, \texttt{RF18}, \texttt{RF19}, \texttt{RF20}, \texttt{RF21}, \texttt{RF22}, \texttt{RNF01}, \texttt{RNF06};
    \item \textbf{Pré-condição}: O discente está matriculado na turma; existe atividade aberta com caso clínico cadastrado;
    \item \textbf{Fluxo Principal}:
    \begin{enumerate}
        \item O estudante abre a atividade e o cronômetro regressivo de 20 minutos é iniciado automaticamente;
        \item O estudante navega pelas abas do prontuário simulado (anamnese, evoluções diárias, prescrições e exames);
        \item O estudante detecta pistas clínicas e seleciona os gatilhos correspondentes no catálogo IHI;
        \item O estudante aponta se houve dano real e classifica a gravidade na escala NCC MERP (E a I);
        \item O estudante acessa a seção de Ferramentas da Qualidade e constrói o Diagrama de Ishikawa 6M em torno do evento adverso;
        \item O estudante pondera os problemas clínicos na Matriz GUT, observando o cálculo dinâmico em tempo real do score $G \times U \times T$;
        \item O estudante estrutura contramedidas na Matriz 5W2H e descreve o ciclo de melhoria PDCA;
        \item O estudante finaliza a submissão dentro do prazo estabelecido.
    \end{enumerate}
    \item \textbf{Fluxos de Exceção}:
    \begin{itemize}
        \item \textbf{Esgotamento do Cronômetro}: Ao atingir 00:00, o sistema emite Toast de encerramento do tempo preconizado pelo IHI e alerta para o fechamento da fase de busca ativa;
        \item \textbf{Prazo da Atividade Expirado}: O sistema bloqueia a submissão se o envio ocorrer após a data limite estipulada pelo docente titular.
    \end{itemize}
\end{itemize}
